%%%%%%%%%%  scan idx 92  |  folio 81  %%%%%%%%%%

\underline{Corolario 3.8.}

\begin{enumerate}[label=(\alph*)]
\item Sean $A$ y $B$ $R$-álgebras tales que
\[
 \operatorname{Tor}_i^R(A,B)=0\qquad(0<i\leq n).
\]
Entonces, para todo haz aditivo
$F\in\operatorname{Ob}(\widetilde{\underline C}_{A\otimes_RB|R})$,
la aplicación canónica
\[
 H^i(A\otimes_RB|R,F)
 \longrightarrow H^i(A|R,F)\oplus H^i(B|R,F)
\]
es un isomorfismo para todo $i\leq n$.

\item Sea $A$ una $R$-álgebra tal que
$\operatorname{Tor}_i^R(A,A)=0$ para todo $i>0$. Si la aplicación
$A\otimes_RA\longrightarrow A$ es un isomorfismo, entonces, para todo
haz aditivo
$F\in\operatorname{Ob}(\widetilde{\underline C}_{A|R})$, se tiene
\[
 H^i(A|R,F)=0
\]
para todo $i$.
\end{enumerate}

\underline{Demostración.}

(a) Sea
$F\in\operatorname{Ob}(\widetilde{\underline C}_{A\otimes_RB|R})$
aditivo. Aplicada a
$R\longrightarrow A\longrightarrow A\otimes_RB$, la Proposición
3.7(c) da la sucesión exacta
\[
 (*)\quad\cdots\longrightarrow
 H^i(A\otimes_RB|A,F_A)
 \longrightarrow H^i(A\otimes_RB|R,F)
 \longrightarrow H^i(A|R,F)\longrightarrow\cdots.
\]
Considérese el diagrama conmutativo
\[
\begin{tikzcd}[column sep=huge,row sep=huge]
H^i(A\otimes_RB|A,F_A) \arrow[r] \arrow[d]
& H^i(A\otimes_RB|R,F) \arrow[d]\\
H^i(B|R,\widetilde A F_A) \arrow[r,"\approx"]
& H^i(B|R,F).
\end{tikzcd}
\]
Aquí
\[
 A:\underline C_{B|R}\longrightarrow
 \underline C_{A\otimes_RB|A},
 \qquad X\longmapsto A\otimes_RX,
\]
y la flecha inferior está inducida por
$\widetilde A F_A\longrightarrow F$, que es un isomorfismo porque $F$ es
aditivo. Por funtorialidad, la conmutatividad del diagrama se reduce al
caso $i=0$.

Si $\operatorname{Tor}_i^R(A,B)=0$ para $0<i\leq n$, entonces la
Proposición 3.7(a) muestra que la flecha vertical izquierda es un
isomorfismo para $i\leq n$. Por consiguiente, tanto los compuestos
\[
 H^i(A\otimes_RB|A,F_A)
 \longrightarrow H^i(A\otimes_RB|R,F)
 \longrightarrow H^i(B|R,F)
\]
como, por simetría,

%%%%%%%%%%  scan idx 93  |  folio 82  %%%%%%%%%%

\[
 H^i(A\otimes_RB|B,F_B)
 \longrightarrow H^i(A\otimes_RB|R,F)
 \longrightarrow H^i(A|R,F)
\]
son isomorfismos para $i\leq n$. Se sigue que $(*)$ se descompone en
sucesiones exactas cortas escindidas
\[
 0\longrightarrow H^i(A\otimes_RB|A,F_A)
 \longrightarrow H^i(A\otimes_RB|R,F)
 \longrightarrow H^i(A|R,F)\longrightarrow0
\]
para $i\leq n$, y por tanto que
\[
 H^i(A\otimes_RB|R,F)
 \longrightarrow H^i(A|R,F)\oplus H^i(B|R,F)
\]
es un isomorfismo en ese rango.

(b) Por (a), la aplicación
\[
 H^i(A\otimes_RA|R,F)
 \longrightarrow H^i(A|R,F)\oplus H^i(A|R,F)
\]
es un isomorfismo para todo $i$. Por otra parte, el isomorfismo supuesto
$A\otimes_RA\longrightarrow A$ induce un isomorfismo
\[
 H^i(A|R,F)\longrightarrow H^i(A\otimes_RA|R,F).
\]
Su compuesto es la aplicación diagonal
\[
 H^i(A|R,F)\longrightarrow
 H^i(A|R,F)\oplus H^i(A|R,F).
\]
Puesto que es un isomorfismo, $H^i(A|R,F)=0$ para todo $i$.
\hfill$\square$

\underline{Corolario 3.9.} Sea $S$ un subconjunto multiplicativo de una
$R$-álgebra $A$. Supongamos que
\[
 F\in\operatorname{Ob}
 (\widetilde{\underline C}_{S^{-1}A|R})
\]
es aditivo y que $F_A$ es también aditivo. Entonces la aplicación canónica
\[
 H^i(S^{-1}A|R,F)\longrightarrow H^i(A|R,F)
\]
es un isomorfismo para todo $i$. Aquí
$F_A\in\operatorname{Ob}
(\widetilde{\underline C}_{S^{-1}A|A})$ está definido por
\[
 F_A(X)=\ker\bigl(F(X)\longrightarrow F(A)\bigr)
\]
para todo
$X\in\operatorname{Ob}(\underline C_{S^{-1}A|A})$.

\underline{Demostración.} Aplicada a
$R\longrightarrow A\longrightarrow S^{-1}A$, la Proposición 3.7(c)
da la sucesión exacta
\[
 \cdots\longrightarrow H^n(S^{-1}A|A,F_A)
 \longrightarrow H^n(S^{-1}A|R,F)
 \longrightarrow H^n(A|R,F)\longrightarrow\cdots.
\]
La aplicación
\[
 S^{-1}A\otimes_AS^{-1}A\longrightarrow S^{-1}A
\]
\begingroup
\sloppy
es un isomorfismo, y $F_A$ es aditivo por hipótesis. Por tanto, el
Corolario 3.8(b) da $H^n(S^{-1}A|A,F_A)=0$ para todo $n$. En consecuencia,
la aplicación central de la sucesión exacta es un isomorfismo para todo
$n$. \hfill$\square$
\par
\endgroup
\par
